\documentclass[11pt,a4paper]{article}
\usepackage[T1]{fontenc}
\usepackage[utf8]{inputenc}
\usepackage[english]{babel}
\usepackage{amsmath,amssymb}
\usepackage[margin=2cm]{geometry}
\setlength{\parindent}{0pt}
\setlength{\emergencystretch}{3em}
\allowdisplaybreaks

\title{About \texttt{sign\_table.tex}}
\author{}
\date{}
\begin{document}
\maketitle
This note explains what the test data in \texttt{sign\_table.tex} is used for. That file is input for the engine in \texttt{main.c}; it is read by the test \texttt{test\_equations\_inequalities.c} in this folder.

\begin{enumerate}
\item \textbf{Purpose.} The file holds 30 true statements in one real or rational variable, produced by the fuzz script \texttt{tools/fuzz/fuzz\_reals.py}. The script's oracle is independent of the engine: it knows the roots of every factor because it chose them, and evaluates the statement exactly on every root and between them. The engine must prove every one.

\item \textbf{What it exercises.} The same rule as the book exercises, on shapes a textbook rarely prints: repeated factors, factors without real roots, irrational roots $\pm\sqrt{d}$ described without radicals (by $x^{2}$ and the sign of $x$), quantifiers $\exists$, negations in front, the sets $\mathbb{Q}$, $\mathbb{R}^{+}$ and $\mathbb{Q}^{+}$.

\item \textbf{Variety.} On the rationals an irrational root is not a value of the variable, so a statement can be true on $\mathbb{Q}$ and false on $\mathbb{R}$; both kinds are present.

\item \textbf{Expected result.} The line \texttt{PASS: 30 of 30 statements proved, each with a written proof}.

\item \textbf{How to run it.} From the top folder of the repository:
\begin{verbatim}
cc -std=c11 -O2 -Wall -Wextra -o /tmp/test_equations_inequalities \
    tests/equations_inequalities/test_equations_inequalities.c \
    && /tmp/test_equations_inequalities
\end{verbatim}
\end{enumerate}
\end{document}
